% Created 2026-03-06 Fri 12:41
% Intended LaTeX compiler: pdflatex
\documentclass[11pt]{article}
\usepackage[utf8]{inputenc}
\usepackage[T1]{fontenc}
\usepackage{graphicx}
\usepackage{longtable}
\usepackage{wrapfig}
\usepackage{rotating}
\usepackage[normalem]{ulem}
\usepackage{amsmath}
\usepackage{amssymb}
\usepackage{capt-of}
\usepackage{hyperref}
\author{The Prologos Project - from Zachary Larson}
\date{2026-03-06}
\title{Effectful Computation on Propagators: The Layered Recovery Principle}
\hypersetup{
 pdfauthor={The Prologos Project - from Zachary Larson},
 pdftitle={Effectful Computation on Propagators: The Layered Recovery Principle},
 pdfkeywords={},
 pdfsubject={},
 pdfcreator={Emacs 30.2 (Org mode 9.7.39)}, 
 pdflang={English}}
\begin{document}

\maketitle
\tableofcontents

\section{The Fundamental Limit}
\label{sec:org3cb0145}

Propagator networks compute monotone fixed points on lattice-valued cells.
Information only grows; values only refine. This is the source of their power
--- convergence is guaranteed, parallelism is safe, and partial information is
first-class.

But IO effects are not monotone. Writing to a file, printing to a console,
opening a network connection --- these are \emph{irreversible observations} that
produce side effects in the world. They cannot be undone, retried, or merged.
The CALM theorem (Hellerstein 2011) formalizes this: non-monotonic operations
require coordination. You cannot achieve confluence for effectful computation
without synchronization.

This creates a tension. Prologos uses propagator networks as its computation
substrate. Session types specify protocols that include IO effects. The
propagator network must somehow handle effects that violate its fundamental
monotonicity assumption.

\textbf{The question is not whether coordination is needed --- it is --- but where
coordination is placed, and how much of the reasoning about effects can remain
monotone.}
\section{The Insight: Layered Recovery}
\label{sec:orgeeb642d}

The solution is not to extend propagators with non-monotone operations.
It is to \emph{layer} a control mechanism on top of the monotone substrate, so
that:

\begin{itemize}
\item \textbf{Reasoning about effects} (what effects exist, in what order, with what
dependencies) happens monotonically in the propagator network
\item \textbf{Executing effects} (performing the actual IO) happens at a control boundary,
sequentially, after the monotone reasoning reaches a fixed point
\end{itemize}

This is the \textbf{Layered Recovery Principle}:

\begin{quote}
Non-monotone behavior is recovered on a monotone substrate by inserting control
layers between phases of monotone computation. The control layers perform the
non-monotone operations (effect execution, negation evaluation, choice
commitment) at precisely the points where coordination is required.
\end{quote}

This principle is not specific to IO effects. It is the general mechanism by
which Prologos extends the propagator substrate to handle any non-monotone
domain.
\section{The Precedent: Logic Engine}
\label{sec:org3934151}

The first application of Layered Recovery in Prologos is the logic engine
(\texttt{rel}), which recovers non-monotone logical operations on the propagator
substrate:

\begin{center}
\begin{tabular}{rll}
Layer & Component & Role\\
\hline
1 & Propagator Network & Monotone constraint propagation (unification)\\
2 & ATMS & Hypothetical reasoning (choice points)\\
3 & Stratification Controller & Non-monotone operations (negation, cut)\\
\end{tabular}
\end{center}

\begin{itemize}
\item \textbf{Layer 1} handles the monotone core: unification constraints propagate through
cells, values refine, the network converges.
\item \textbf{Layer 2} (ATMS) manages alternative worldviews. Each choice point creates an
assumption; the ATMS tracks which combinations of assumptions are consistent
(and which are nogoods). This is still monotone --- adding nogoods only
shrinks the set of valid worldviews.
\item \textbf{Layer 3} is the stratification barrier. Negation (\texttt{\textbackslash{}+} in Prolog) is
non-monotone: "P is not provable" can become false if P is later derived.
Stratification resolves this by dividing the program into strata, computing
each stratum to a fixed point before evaluating negation that references
lower strata. The barrier between strata is the coordination point.
\end{itemize}

The key observation: Layers 1 and 2 are entirely monotone and live within the
propagator network. Layer 3 is the \emph{minimal} non-monotone extension, placed
at exactly the boundary required by the CALM theorem. The logic engine doesn't
abandon the propagator model --- it \emph{layers} control on top of it.

See: \texttt{docs/tracking/2026-02-24\_LOGIC\_ENGINE\_DESIGN.org} §5 (Three-Layer
Architecture), Appendix B (Stratification).
\section{The Application: Session-Derived Effect Ordering}
\label{sec:org4e5fb17}

The second application of Layered Recovery handles IO effects. The insight that
unlocks it: \emph{session types are causal clocks}. Each position in a session type's
continuation chain is a causally-ordered event. The ordering is not imposed
externally --- it is \textbf{intrinsic} to the session type's structure.
\subsection{Session Types as Causal Clocks}
\label{sec:org73b94e7}

A session type \texttt{!String . ?Int . end} encodes three positions:

\begin{verbatim}
Position 0: send a String    (happens first)
Position 1: recv an Int      (happens second)
Position 2: end              (happens last)
\end{verbatim}

Session advancement --- the monotone process of advancing a session cell from
one protocol state to the next --- is a \emph{Lamport clock}. Each advancement ticks
the clock forward. For multiple channels, the per-channel clocks form a \emph{vector
clock}: operations on different channels are concurrent (incomparable in the
partial order), while operations on the same channel are totally ordered.

This is the same causal structure that distributed systems use for ordering
events, but derived from the \emph{type system} rather than assigned at runtime.
Session fidelity (a theorem of the session type system) guarantees that effects
respect this ordering.
\subsection{The Galois Connection}
\label{sec:org880c6c3}

The formal bridge between sessions and effects is a Galois connection:

\begin{verbatim}
 (alpha, gamma)
alpha : Session -> EffectPosition      (extract causal position from session state)
gamma : EffectPosition -> Session      (reconstruct remaining protocol from position)
\end{verbatim}

Both functions are monotone. The adjunction guarantees soundness: session
fidelity implies effect ordering correctness. At a deeper level, both the
session lattice and the effect position lattice are \emph{effect quantales} ---
lattice-ordered monoids with sequential composition --- and the Galois
connection is a quantale morphism preserving both ordering and composition.

This formalizes why session types can serve as the foundation for effect
ordering: the algebraic structure is shared. The connection is not accidental;
it reflects a deep mathematical relationship between protocols and effects.

See: \texttt{docs/tracking/2026-03-06\_SESSION\_TYPES\_AS\_EFFECT\_ORDERING.org} §4 (Galois
Connection), §6 (Effect Quantale Connection).
\subsection{Cross-Channel Data Dependencies}
\label{sec:org6cdfc43}

Session types provide per-channel ordering but not cross-channel ordering.
When a value received on channel \texttt{a} flows to a send on channel \texttt{b}, there
is a data dependency that the session types alone don't capture.

The resolution: compute the \textbf{transitive closure} of the union of session
ordering edges and data-flow edges. Transitive closure is a monotone
operation on sets of edges --- adding an edge never removes existing ordering
relationships. This means a propagator can incrementally compute the complete
partial order:

\begin{verbatim}
Session edges:    (a0 < a1), (b0 < b1)           (from session types)
Data-flow edge:   (a0 < b0)                       (x flows from recv a to send b)
Transitive close: (a0 < a1), (b0 < b1), (a0 < b0), (a0 < b1)
\end{verbatim}

The transitive closure is a monotone fixed point. It lives natively in the
propagator network. No external scheduling is needed.
\subsection{The Five-Layer Architecture}
\label{sec:org4b1f51e}

The full architecture for effectful computation on propagators:

\begin{verbatim}
Layer 1: Session Advancement              [exists -- session-runtime.rkt]
    Session cells advance monotonically through session types.
    Each advancement is a causal clock tick for that channel.

Layer 2: Data-Flow Analysis               [NEW]
    Analyze variable bindings to derive cross-channel ordering edges.
    recv x a -> send [f x] b  ==>  ordering edge (a_pos < b_pos)

Layer 3: Transitive Closure               [NEW]
    Compute the complete partial order by propagating ordering edges
    to their transitive closure. Monotone fixed-point computation.
    Concurrent effects remain unordered (least commitment).

Layer 4: ATMS Branching                   [exists -- atms.rkt]
    Maintain per-branch ordering hypotheses for proc-case.
    Resolve when branch is chosen; discard inconsistent worldviews.

Layer 5: Effect Handler                   [NEW]
    Read the resolved partial order. Execute effects in any valid
    linearization. This is the only non-monotone step (actual IO).
\end{verbatim}

Layers 1--4 are monotone. Layer 5 is the control boundary --- the precise
point where the CALM theorem requires coordination.
\section{The \texttt{rel\textasciitilde{}/\textasciitilde{}proc} Parallel}
\label{sec:orga17b511}

The pattern reveals a deep structural parallel between Prologos's two
computational engines:

\begin{center}
\begin{tabular}{lll}
Dimension & \texttt{rel} (Logic Engine) & \texttt{proc} (Session Runtime)\\
\hline
\textbf{Variables} & Logic variables (unification cells) & Session cells (protocol state)\\
\textbf{Constraints} & Equality constraints (type/term) & Ordering constraints (causal)\\
\textbf{Propagation} & Unification propagators & Session advancement + ordering propagators\\
\textbf{Fixed point} & Constraint closure (all equalities) & Ordering closure (transitive causal edges)\\
\textbf{Hypotheses} & ATMS worldviews (choice points) & ATMS worldviews (branch alternatives)\\
\textbf{Resolution} & Evidence selects worldview & Received value selects worldview\\
\textbf{Execution} & Proof terms (constructive evidence) & Effects (IO operations in causal order)\\
\textbf{Control layer} & Stratification (negation evaluation) & Effect handler (IO execution)\\
\end{tabular}
\end{center}

Both engines share the three-layer architecture: Propagator Network (monotone
fixed points) + ATMS (hypothetical reasoning) + Control Layer (non-monotone
operations). The Layered Recovery Principle is the \emph{common theory} underlying
both.

This parallel is not coincidental. Both \texttt{rel} and \texttt{proc} face the same
fundamental challenge: how to perform non-monotone operations (negation for
\texttt{rel}; effects for \texttt{proc}) on a monotone substrate (propagator network). Both
solve it the same way: reason monotonically about the non-monotone domain, then
execute the non-monotone operations at a control barrier after the monotone
reasoning converges.
\section{What This Means: Designing for the Principle}
\label{sec:orgcd6a137}

The Layered Recovery Principle is a design methodology, not just an observation.
When extending Prologos to handle a new domain that involves non-monotone
operations:

\begin{enumerate}
\item \textbf{Identify the monotone core}: What can be reasoned about convergently?
For effects, it's ordering relationships. For logic, it's unification
constraints. For type inference, it's subtyping relationships.

\item \textbf{Encode the monotone core in the propagator network}: Define lattices,
cells, and propagators for the convergent reasoning. This gets parallelism,
incrementality, and partial-information handling for free.

\item \textbf{Identify the non-monotone boundary}: What operations require
coordination? For effects, it's execution. For logic, it's negation. This
is where the CALM theorem demands synchronization.

\item \textbf{Insert a control layer at the boundary}: The control layer reads the
fixed point of the monotone core and performs the non-monotone operations
in a coordinated fashion.

\item \textbf{Use the ATMS for hypothetical reasoning}: If the domain involves
alternatives (choice points, branches, speculative evaluation), the ATMS
manages the hypotheses monotonically.
\end{enumerate}

This methodology has now been applied successfully to two domains (logic,
effects) and the consistent results suggest it will generalize to others.
\section{The Architecture Decision: A + D}
\label{sec:org24f402c}

Two architectures for effect ordering:

\begin{itemize}
\item \textbf{Architecture A (Stratified Barriers)}: Effects execute in AST walk order.
Sequential, simple, correct by construction. No propagator involvement in
ordering. This is the current implementation for \texttt{main} and top-level
REPL expressions.

\item \textbf{Architecture D (Session-Derived Ordering)}: Effects are ordered by session
type structure via the Galois connection. Cross-channel dependencies resolved
by transitive closure. Branches handled by ATMS. The ordering is a theorem
from session fidelity, not a design choice.
\end{itemize}

The recommended architecture is \textbf{A + D}:

\begin{itemize}
\item \textbf{D} for session-typed IO (the structured, common case). Ordering is derived
from session types, verified by the type checker, and compositional across
protocol phases. Concurrent effects on independent channels are correctly
unordered (least commitment).

\item \textbf{A} for unsessioned IO (the REPL, top-level eval, \texttt{main} without explicit
sessions). Walk-based ordering is simple and sufficient.
\end{itemize}

Architecture A is, in fact, a \textbf{degenerate case} of Architecture D: \texttt{main} with
sequential IO and no explicit session is equivalent to a process with a single
implicit session. As \texttt{main} gains explicit session structure, D's richer
ordering semantics apply automatically.

Architecture B (Timestamped Effect Cells) is \textbf{subsumed}: for session-typed
effects, D provides strictly more information (intrinsic, verified ordering);
for unsessioned effects, A is simpler. B adds nothing that A + D don't provide.
\section{Relationship to Other Principles}
\label{sec:org9f2b300}

\begin{center}
\begin{tabular}{ll}
Document & Relationship\\
\hline
\texttt{DESIGN\_PRINCIPLES.org} & "Correctness Through Types" --- effect ordering is derived from session types (correctness is a theorem). "Layered Architecture" --- Layers 5/6 (session/process) are where effect ordering operates.\\
\texttt{PROTOCOLS\_AS\_TYPES.org} & Session type composition (protocol libraries) composes effect orderings automatically. Capability composition tracks alongside protocol and effect composition.\\
\texttt{LANGUAGE\_VISION.org} & "Protocol correctness" extends to effect correctness. "Session types as protocol specifications" gains causal clock interpretation.\\
\texttt{CAPABILITY\_SECURITY.md} & \texttt{SysCap} as powerbox for \texttt{main} (Architecture A fallback). Capability requirements compose with session phases and their effect orderings.\\
\texttt{RELATIONAL\_LANGUAGE\_VISION.org} & The logic engine (\texttt{rel}) is the first application of Layered Recovery. The structural parallel between \texttt{rel} and \texttt{proc} confirms the principle's generality.\\
IO Implementation PIR & §4a ("Propagators Cannot Order Side Effects") is the discovery that motivated this principle. The PIR documents the empirical observation; this document provides the theoretical resolution.\\
\end{tabular}
\end{center}
\section{Implications for the Language}
\label{sec:org4c85e35}

\subsection{Effect Ordering Is a Type-System Property}
\label{sec:org1c3c54c}

With Architecture D, the ordering of IO effects is not a scheduler decision
or a compiler implementation detail. It is a \textbf{property of the type system}.
The session type determines the ordering; the type checker verifies it; the
compiler derives it via the Galois connection. Users can reason about effect
ordering by reading the session type --- the same way they reason about protocol
correctness.
\subsection{Concurrent Effects Fall Out Naturally}
\label{sec:org3258c99}

For multi-channel processes, effects on independent channels are concurrent ---
the partial order leaves them unordered. This means the runtime can execute
them in parallel or in any interleaving, without violating correctness. The
session type system guarantees that these effects are on independent resources
(different channels, verified by linearity). Concurrency is not opt-in; it's
the default for independent effects.
\subsection{Static Deadlock Detection}
\label{sec:org504549f}

The transitive closure computation for cross-channel data dependencies can
detect cycles (\texttt{A < B < A}), which indicate deadlocks. This turns runtime
deadlocks into compile-time errors. The ordering analysis serves double duty:
it provides effect ordering \emph{and} deadlock detection for data-dependent
cross-channel interactions.
\subsection{The CALM Optimality}
\label{sec:org24a4831}

The architecture achieves the minimum coordination required by the CALM theorem.
All reasoning about effects (ordering, dependencies, alternatives) is monotone
and coordination-free. The only coordination point is the final effect execution
barrier (Layer 5). This is provably optimal --- no architecture can eliminate
this barrier without violating the CALM theorem.
\section{Theoretical Foundations}
\label{sec:orgda49677}

\subsection{CALM Theorem (Hellerstein 2011)}
\label{sec:org7493866}

The Consistency And Logical Monotonicity theorem: a distributed program has a
consistent, coordination-free implementation if and only if it is monotonic.
IO effects are non-monotonic, so coordination is required. The Layered Recovery
Principle minimizes this coordination to a single barrier.
\subsection{Propagator Networks (Radul \& Sussman 2009)}
\label{sec:orgdb9b339}

Monotone fixed-point computation on lattice-valued cells. The computational
substrate for Prologos. Propagator designers explicitly excluded non-commutative
effects from the model --- the Layered Recovery Principle is our answer to this
exclusion.
\subsection{Session Types (Caires \& Pfenning 2010, Wadler 2012)}
\label{sec:org74eeca2}

Session types as linear logic propositions via Curry-Howard. The continuation
structure of a session type encodes the causal ordering of communication events.
Session fidelity guarantees that implementations respect this ordering.
\subsection{Galois Connections}
\label{sec:orga9a9a9a}

An adjunction \texttt{(\textbackslash{}alpha, \textbackslash{}gamma)} between ordered domains. Used throughout abstract
interpretation (Cousot \& Cousot) for soundly approximating one domain in terms
of another. Our use connects the session lattice to the effect position lattice.
\subsection{Vector Clocks (Lamport 1978, Fidge/Mattern 1988)}
\label{sec:orga87232b}

Per-process Lamport clocks composed into vectors for partial ordering of events
in distributed systems. Per-channel session depths composed into vectors for
partial ordering of effects across channels. The mathematical structure is
identical.
\subsection{Partial Order Planning (Weld 1994)}
\label{sec:orgc510ee4}

Least-commitment planning with causal links and ordering constraints. Effect
ordering via transitive closure is isomorphic to POP: data-flow edges are
causal links, session ordering is initial constraints, \texttt{proc-case} is plan
alternatives, ATMS is threat postponement.
\subsection{Effect Quantales (Katsumata 2014, Gordon 2021)}
\label{sec:org7445308}

Lattice-ordered monoids with sequential composition that distribute over joins.
Both session types and effect positions form effect quantales. The Galois
connection is a quantale morphism preserving this algebraic structure.
\subsection{Category-Graded Monads (Katsumata, Orchard et al.)}
\label{sec:org3300b4d}

Unifies graded monads (effect systems) and parameterised monads (session types)
into a single categorical framework. The Galois connection between sessions
and effects is, categorically, a functor between indexing categories.
\section{Inspirations}
\label{sec:orgf70071c}

\begin{itemize}
\item \textbf{Hellerstein} (CALM theorem) --- The theoretical foundation for why
coordination is necessary and where it must be placed.
\item \textbf{Radul \& Sussman} (Propagators) --- The computational substrate. Their
explicit exclusion of non-commutative effects is the challenge we resolve.
\item \textbf{Caires \& Pfenning} (Session types as linear logic) --- The Curry-Howard
foundation that makes session types causal clocks.
\item \textbf{Lamport} (Logical clocks) --- Session advancement IS a Lamport clock.
The connection is direct, not analogical.
\item \textbf{de Kleer} (ATMS) --- Hypothetical reasoning for branching alternatives.
The same ATMS serves both \texttt{rel} and \texttt{proc}.
\item \textbf{Weld} (Partial Order Planning) --- The least-commitment principle for
ordering: don't impose unnecessary ordering on concurrent effects.
\item \textbf{Atkey} (Parameterised monads) --- Session types as pre/post state indices.
Sequential composition = bind matching postcondition to precondition.
\item \textbf{The Prologos Logic Engine} --- The first application of Layered Recovery,
which established the three-layer architecture and demonstrated that
non-monotone operations can be recovered on a monotone substrate.
\end{itemize}
\end{document}
